\paragraph{Periodic lifts: an independent arithmetic audit.}
Put $x=N/p=t+f$ with $t\in\mathbb Z_{\ge0}$ and $0\le f<1$.
In the notation of the six-shell argument, translation by $t$ changes
\[
 a=8t+a_0,\quad M=11t+M_0,\quad
 \lambda_u=ut+\lambda_{u,0},\quad K=13t+K_0.
\]
It leaves $b/p$, $d/p$, $h$, $e$, $\delta$, and the minimal-carry support
unchanged. Thus the same eight rows of the chamber table apply, with these
integer translations. The high-digit polynomial $U(V)$ changes, but it is
still independent of the residue $k$ on the minimal-carry support; that
independence, rather than its explicit coefficients, is the required property.
The degree condition becomes exactly
\[
 13N+r<(13t+K_0)p.
\]
The admissible shift can uniformly be chosen as
$\eta=p\lceil2N/p\rceil$. For $t\ge1$ it lies in $[2N,3N]$;
for $t=0$ it is the original choice and lies in $[2N,9N]$.
All original factorial arguments, polynomial degrees, and harmonic indices
are still below $p^2$ under $p^2>50N$.
The omitted internal chamber walls $f=5/11,10/11$ still force $p=11$,
which is incompatible with the cutoff. Hence no new obstruction appears
in translating the existing local factorization and Newton-sum argument.
This audit must be read with that argument, not as a replacement for its
integer-valued-polynomial justification.

\paragraph{Density of the lifted prime intervals.}
For
\[
 \mathcal B=\{(2,8,1),(3,9,2),(4,10,3),
                 (5,8,3),(6,9,4),(12,9,8)\},
\]
write $(k,a,b)$ for an element. The additional intervals, beyond the six
original intervals, are
\[
 \frac{13N+r}{13t+k}<p<\frac{aN}{at+b},\qquad t=1,2,\ldots .
\]
They are disjoint because their reciprocals lie in the disjoint translated
base intervals. If $H_{N,r}^{\mathrm{lift}}$ is the product of eligible primes
in these intervals, the prime number theorem gives
\[
 \lim_{N\to\infty}\frac{\log H_{N,r}^{\mathrm{lift}}}{N}
 =g:=\sum_{(k,a,b)\in\mathcal B}\sum_{t=1}^{\infty}
       \left(\frac{a}{at+b}-\frac{13}{13t+k}\right).
\]
To justify the infinite sum, first retain finitely many $t$. The column
shift changes each endpoint only by a bounded amount. The discarded
intervals for $t\ge T$ contain only primes $p<N/T$, so their total
logarithmic weight is at most $\vartheta(N/T)$; after division by $N$,
the limiting upper bound is $1/T$. The excluded primes
$p\le\sqrt{50N}$ contribute $o(N)$. This proves the stated limit without
requiring a prime number theorem uniform in $t$.
Equivalently,
\[
 g=\sum_{(k,a,b)\in\mathcal B}
       \bigl(\psi(1+k/13)-\psi(1+b/a)\bigr)
  =0.10868236921910157039\ldots.
\]

\paragraph{An exact rational enclosure.}
Each summand is
\[
 f_{k,a,b}(t)=\frac{ak-13b}{(at+b)(13t+k)}>0.
\]
Put $M=100000$ and $S=10^{24}$. Integer division gives
\[
 A:=\sum_{(k,a,b)\in\mathcal B}\sum_{t=1}^{M}
          \lfloor S f_{k,a,b}(t)\rfloor
   =108681309400407200714213.
\]
Since $f(t)=(\beta-\alpha)/((t+\alpha)(t+\beta))$ with
$\alpha=b/a<\beta=k/13$, comparison with decreasing reciprocal squares gives
\[
 \frac{ak-13b}{a(13M+13+k)}
 \le \sum_{t=M+1}^{\infty}f_{k,a,b}(t)
 \le \frac{ak-13b}{13aM}.
\]
Consequently
\[
 \frac AS+\sum_{\mathcal B}\frac{ak-13b}{a(13M+13+k)}
 \le g\le
 \frac{A+6M}{S}+\sum_{\mathcal B}\frac{ak-13b}{13aM}.
\]
The accompanying integer-arithmetic script verifies this enclosure:
\[
 0.1086823692136766<g<0.1086823692294671.
\]
Using the previously certified carry-density lower bound
\[
 \gamma\ge\frac{159010385641752609841}{13110552000000000000},
\]
set $c_0=25-233/120-\gamma_{\rm lower}$ and let $g_-$ denote the displayed
rational lower endpoint. The existing analytic enclosures imply
\[
 \sigma\le38.168732909647+c_0-g_-,\qquad
 \tau\ge24.467842131921-c_0+g_->0.
\]
Exact rational cross-multiplication then gives
\[
 1+\frac{38.168732909647+c_0-g_-}
          {24.467842131921-c_0+g_-}
   <4.589897.
\]
This last implication uses the uniform integer normalization furnished by
the translated shell argument and the analytic/nonvanishing inputs.
It does not follow merely from the finite divisibility checks below.

\paragraph{Why the column shift matters even after translation.}
At $N=3000$, $p=907$, the active lift has $t=3$, $k=4$ and
$(13t+k)p=43\cdot907=39001$.
Thus $13N+r<(13t+k)p$ holds only for $r=0$.
For the pair normalized using the original six-shell divisor, its two
$p$-valuations are $(1,3)$ in column zero, but $(0,2)$ in columns one and
two. The unshifted periodic guess therefore fails, whereas the shifted
inequality predicts precisely the observed distinction.
The same phenomenon occurs at $N=6000$, $p=1013$:
$(13\cdot5+12)1013=78001$.
